% ================================================================== COURSE OVERVIEW
\chapter*{Course Overview}
\addcontentsline{toc}{chapter}{Course Overview}
\markboth{Course Overview}{}

\section*{Course description}
Business Process Analytics trains you to work as a \emph{process and operations analyst}: the person who turns process
data into decisions about stock, staff, capacity, routes and processes, and who can show what those decisions are
worth in francs. You already know how to fit models. This course is about everything around the model that
decides whether it creates value: what decision it serves, what an error costs, which constraints bind, how
uncertain the answer is, and whether the change actually worked.

The course follows one thread. Each week's output is the next week's input, as it would be in a real analytics team.

\begin{figure}[H]
\centering
\begin{tikzpicture}[node distance=3mm,
  stage/.style={draw=none,fill=#1,text=white,font=\sffamily\bfseries\footnotesize,minimum height=12mm,
    text width=17mm,align=center,rounded corners=1.5pt},
  lab/.style={font=\sffamily\scriptsize,text=slate,align=center,text width=19mm}]
  \node[stage=navy] (w1) {W1\\Measure};
  \node[stage=navy,right=of w1] (w2) {W2\\Diagnose};
  \node[stage=teal,right=of w2] (w3) {W3\\Forecast};
  \node[stage=teal,right=of w3] (w4) {W4\\Stock policy};
  \node[stage=orange,right=of w4] (w5) {W5\\Optimise};
  \node[stage=orange,right=of w5] (w6) {W6\\Capacity \& risk};
  \node[stage=brick,right=of w6] (w7) {W7\\Did it work?};
  \foreach \a/\b in {w1/w2,w2/w3,w3/w4,w4/w5,w5/w6,w6/w7} \draw[-{Latex[length=2mm]},slate,thick] (\a)--(\b);
  \node[lab,below=2mm of w1] {clean data,\\KPIs, ABC--XYZ};
  \node[lab,below=2mm of w2] {control charts,\\process maps, alerts};
  \node[lab,below=2mm of w3] {point \& quantile\\forecasts, errors};
  \node[lab,below=2mm of w4] {safety stock,\\newsvendor, bullwhip};
  \node[lab,below=2mm of w5] {LP/MIP plans,\\shadow prices, routes};
  \node[lab,below=2mm of w6] {staffing to SLA,\\simulation, VSS};
  \node[lab,below=2mm of w7] {A/B, DiD,\\AI guardrails};
  \draw[decorate,decoration={brace,amplitude=5pt,raise=3pt},teal,thick] (w3.north west)--(w6.north east)
    node[midway,above=9pt,font=\sffamily\bfseries\footnotesize,text=teal]{predict-then-optimise: the forecast's errors set the decision's buffers};
\end{tikzpicture}
\caption{The decision thread. Colours mark the three questions: what is happening (navy), what will happen and what to do about it (teal, orange), and did it work (red).}
\end{figure}

\section*{Learning outcomes}
By the end of the course a successful student can:
\begin{enumerate}[label=\textbf{LO\arabic*.},leftmargin=3.2em]
\item frame an operational decision: objective, decision variables, constraints, KPI and the cost of each kind of error;
\item build a trustworthy analytical dataset from raw operational records with SQL and Python, with every cleaning rule documented and counted;
\item diagnose a process with control charts, Little's Law, process mining and cost-based anomaly alerts;
\item produce point and probabilistic forecasts and validate them honestly against simple baselines;
\item convert forecasts and their errors into inventory policies and test those policies by simulation;
\item formulate and solve linear, integer and routing models, and interpret shadow prices and sensitivity;
\item size capacity against a service-level agreement with queueing theory and discrete-event simulation, and plan under uncertainty with stochastic optimisation;
\item estimate the causal effect of an operational change and deploy AI assistance with explicit guardrails;
\item recommend an action quantified in XAF, with its uncertainty, and defend it orally.
\end{enumerate}

\section*{Seven-week map}
\begin{center}\small
\begin{tabularx}{\linewidth}{@{}c L{3.0cm} Y L{2.5cm} L{1.7cm}@{}}
\toprule
\textbf{Wk} & \textbf{Theme} & \textbf{Core methods and lab} & \textbf{Real data} & \textbf{LOs}\\\midrule
1 & Operational data, KPIs and SQL & KPI design, data-quality rules as code, DuckDB SQL with window functions, ABC--XYZ segmentation & Online Retail II (1.07\,M lines) & 1, 2, 9\\
2 & Process diagnostics & Common vs special cause, $p$-charts, Phase I/II limits, Little's Law, process discovery and conformance, alerts tuned by cost & Permit process log; NAB machine temperature & 3\\
3 & Forecasting for decisions & Rolling-origin validation, WAPE/MASE/bias, ETS, Croston/ADIDA, LightGBM benchmark, conformal intervals & Online Retail weekly SKU demand & 4\\
4 & Inventory policy from forecasts & EOQ, safety stock from forecast error, $(s,Q)$ and $(R,S)$ simulation, newsvendor with quantiles, bullwhip & Lab 3 forecasts; beer game & 5\\
5 & Optimisation & LP and duality, MIP modelling, facility location, rostering, vehicle routing & Cameroon city network; Douala deliveries & 1, 6\\
6 & Capacity, queues and uncertainty & Erlang C staffing to an SLA, SimPy simulation with replications, two-stage stochastic programming, VSS and EVPI & Online Retail order arrivals & 7\\
7 & Did it work? & A/B testing with SRM checks, regression and doubly robust adjustment, DiD and event studies, LLM guardrails & Hillstrom e-mail test; organ-donation policy & 8, 9\\
\midrule
Exam wk & Capstone defence & 12-minute presentation, 8-minute questioning on a decision you recommend & Your choice & 1--9\\
\bottomrule
\end{tabularx}
\end{center}

\section*{Anatomy of a 4-hour session}
Every session follows the same rhythm, so no time is lost on set-up.
\begin{center}
\begin{tikzpicture}[x=0.06cm]
  \foreach \s/\e/\c/\t in {0/20/navy/Briefing,20/70/teal/Concepts,70/160/orange/Guided lab,160/220/brick/Challenge,220/240/slate/Debrief}{
    \fill[\c] (\s,0) rectangle (\e,0.75);
    \node[font=\sffamily\scriptsize\bfseries,text=white] at ({(\s+\e)/2},0.375) {\t};}
  \foreach \x/\l in {0/0:00,20/0:20,70/1:10,160/2:40,220/3:40,240/4:00}
    \node[font=\sffamily\scriptsize,text=slate,below] at (\x,0) {\l};
\end{tikzpicture}
\end{center}
About 170 of the 240 minutes are hands-on. The debrief is compulsory: each student states, in two sentences, what the
numbers mean for a manager and what decision they support. Ten minutes of break are taken inside the guided lab.

\section*{Assessment}
\begin{center}\small
\begin{tabularx}{\linewidth}{@{}L{4.1cm} L{2.0cm} Y r@{}}
\toprule
\textbf{Component} & \textbf{When} & \textbf{Content} & \textbf{Weight}\\\midrule
CA1 (practical, 2 h) & after Week 2 & Data audit, KPIs and a process diagnosis on an unseen dataset & 10\%\\
CA2 (practical, 2 h) & after Week 4 & Forecast a set of items and set their stock policy; cost the result & 15\%\\
CA3 (practical, 2 h) & after Week 6 & Formulate and solve an allocation model; test it under uncertainty by simulation & 15\%\\
Capstone checkpoints & Weeks 3 and 5 & Proposal and data audit (W3); baseline model and decision prototype (W5) & 5\%\\
Capstone report and code & exam week & A decision recommendation with quantified impact, reproducible repository & 25\%\\
Capstone defence & exam week & 12-minute presentation, 8 minutes of questions on choices and trade-offs & 20\%\\
Final examination & end of course & Short theory and applied problems (formula sheet allowed) & 10\%\\\midrule
\textbf{Total} & & & \textbf{100\%}\\\bottomrule
\end{tabularx}
\end{center}

\paragraph{Learning outcomes and assessment} Every outcome is assessed at least twice, once in a timed practical
and once in the capstone.
\begin{center}\small\sffamily
\begin{tabular}{@{}lccccccccc@{}}
\toprule
 & LO1 & LO2 & LO3 & LO4 & LO5 & LO6 & LO7 & LO8 & LO9\\\midrule
CA1 & \textbullet & \textbullet & \textbullet & & & & & & \textbullet\\
CA2 & \textbullet & & & \textbullet & \textbullet & & & & \textbullet\\
CA3 & \textbullet & & & & & \textbullet & \textbullet & & \textbullet\\
Capstone (report) & \textbullet & \textbullet & $\circ$ & $\circ$ & $\circ$ & $\circ$ & $\circ$ & \textbullet & \textbullet\\
Defence & \textbullet & & & & & & & \textbullet & \textbullet\\
Final exam & \textbullet & & \textbullet & \textbullet & \textbullet & \textbullet & \textbullet & \textbullet & \\\bottomrule
\multicolumn{10}{@{}l}{\footnotesize \textbullet\ always assessed \quad $\circ$\ assessed when the project uses that method (each project must use at least three)}
\end{tabular}
\end{center}

\paragraph{Lab marking rubric (CA1--CA3)}
\begin{center}\small
\begin{tabularx}{\linewidth}{@{}L{3.6cm}Yr@{}}
\toprule
\textbf{Criterion} & \textbf{What earns full marks} & \textbf{Share}\\\midrule
Correctness & Runs top to bottom on a fresh kernel; no leakage; results reproduce & 30\%\\
Decision framing & States the decision, the cost of each error and the constraint that binds & 20\%\\
Method choice & Justified, compared with a simple baseline, uncertainty shown & 20\%\\
So what & A manager-ready paragraph with numbers, units, XAF impact and a recommended action & 20\%\\
Code and communication & Clear names, comments, labelled axes, no dead code & 10\%\\\bottomrule
\end{tabularx}
\end{center}

\begin{pitfall}{AI assistants and academic integrity}
You may use AI coding assistants. You must (i) be able to explain every line you submit, (ii) list the prompts you
used in an appendix, and (iii) never submit output you have not run. In the defence, a line you cannot explain is
treated as copied. Never paste real personal or confidential organisational data into an external AI service.
\end{pitfall}

\section*{How to read these notes}
Coloured boxes help you find what you need. Code appears in blue-tinted listings; the grey block underneath is
the output the code printed when these notes were built. If your numbers differ beyond rounding, check your
package versions against \texttt{requirements.txt} first.

\begin{center}\small\sffamily
\begin{tabular}{@{}l l@{}}
\tcbox[colback=navy,colframe=navy,coltext=white,size=small,on line]{Definition} & precise statement of a concept; learn these\\
\tcbox[colback=teal,colframe=teal,coltext=white,size=small,on line]{Key formula} & results you must be able to use and derive\\
\tcbox[colback=green,colframe=green,coltext=white,size=small,on line]{Worked example} & a fully calculated numerical example\\
\tcbox[colback=orange,colframe=orange,coltext=white,size=small,on line]{Practical lab} & the guided coding session; type the code, do not paste\\
\tcbox[colback=brick,colframe=brick,coltext=white,size=small,on line]{Common pitfall} & mistakes that cost marks, and money in practice\\
\tcbox[colback=navy!85,colframe=navy!85,coltext=white,size=small,on line]{In practice} & a documented case from a real organisation\\
\end{tabular}
\end{center}

% ================================================================== GETTING STARTED
\chapter*{Getting Started: Tools, Data and the Running Case}
\addcontentsline{toc}{chapter}{Getting Started: Tools, Data and the Running Case}
\markboth{Getting Started: Tools, Data and the Running Case}{}

\section*{A free, low-setup, professional stack}
Python is the single language of the course, so you install one environment and use it all semester. Nothing
requires a paid licence or a Windows machine.

\begin{center}\small
\begin{tabularx}{\linewidth}{@{}L{3.0cm}L{4.4cm}Y@{}}
\toprule
\textbf{Need} & \textbf{Tool} & \textbf{Why}\\\midrule
Run code & Google Colab (primary); local JupyterLab with \texttt{uv} or Miniforge (fallback) & Colab needs no installation and works on weak laptops; the local fallback covers unreliable internet.\\
Version control, submission & GitHub (free) & Industry standard; your capstone repository is what you defend.\\
Reports & Jupyter or Quarto & Free and reproducible; export to HTML or PDF.\\
SQL & DuckDB & No server; queries CSV and Parquet files directly; fast on a million rows.\\
Dashboards & Looker Studio (browser) & Free, shareable, no Windows dependency. A Plotly or Quarto dashboard is an alternative.\\
Forecasting & statsforecast, LightGBM & Fast classical models with honest cross-validation; one ML benchmark.\\
Process mining & PM4Py (needs Graphviz for drawings) & Academic-grade discovery and conformance.\\
Optimisation & PuLP + CBC, OR-Tools & Open-source solvers that handle every model in the course; Google Sheets or LibreOffice Solver for spreadsheet work.\\
Simulation & NumPy, SimPy & Monte Carlo and discrete-event simulation in a few dozen lines.\\
Causal inference & statsmodels, scikit-learn & Regression, propensity models, clustered standard errors.\\\bottomrule
\end{tabularx}
\end{center}

\begin{lstlisting}[style=sh,caption={Setting up (local). In Colab, run the same \texttt{pip} line with a leading \texttt{!}.}]
git clone https://github.com/<your-account>/bpa-yibs.git && cd bpa-yibs
uv venv && source .venv/bin/activate        # Windows: .venv\Scripts\activate
uv pip install -r requirements.txt          # or: pip install -r requirements.txt
python datagen.py                           # Mboa Fresh Market teaching case -> ./data
python get_data.py                          # real public datasets -> ./real
jupyter lab
\end{lstlisting}

\begin{pitfall}{Unpinned versions break labs}
While these notes were being tested, two library releases broke working code. PuLP 4 changed how variables are
created, and an OR-Tools build failed to load when a newer \texttt{highspy} had been imported first (both bundle the
HiGHS solver). Both are pinned in \texttt{requirements.txt}. Put your own \texttt{pip freeze > requirements.txt} in
every submission: ``it worked on my laptop'' is not evidence.
\end{pitfall}

\paragraph{Practical tips} The instructor provides one starter repository with the install cell, the data scripts
and the lab skeletons, so week one is not spent debugging environments. All datasets are also copied to a shared
drive and a USB stick before the course starts, in case the connection fails during a lab.

\section*{Two kinds of data, on purpose}
Every week pairs \textbf{synthetic} data, where the true answer is known, with \textbf{real} data, where it is not
and the data are messy in ways nobody planned.

\paragraph{The teaching case: Mboa Fresh Market} A fictional supermarket chain with six stores in Yaoundé, Douala and
Bafoussam and twelve fast-moving products. \texttt{datagen.py} (Appendix A) writes the files below with a fixed seed.
Use Mboa to check that a method recovers a truth you know, for example the $+8\%$ loyalty-app effect or the $+3$-point
training effect. Do not read the generator to find the answers before Week 7.

\begin{center}\small
\begin{tabularx}{\linewidth}{@{}L{2.6cm}L{2.2cm}Y@{}}
\toprule
\textbf{File} & \textbf{Size} & \textbf{Content and deliberate problems}\\\midrule
\texttt{sales.csv} & 52\,672 rows & Daily units and price by store and product, 2023--2024; duplicates, missing and negative units, keying errors ($\times40$), mixed-case store codes. Douala stores gain $+8\%$ from 1 July 2024.\\
\texttt{products.csv} & 12 rows & Cost, price, category, lead time, order cost and shelf life.\\
\texttt{stores.csv} & 6 rows & City and floor area.\\
\texttt{event\_log.csv} & 9\,193 events & 1\,500 student-registration cases with rework and order violations.\\
\texttt{clinic.csv} & 3\,021 rows & 30 days of walk-in arrivals and service times (queueing exercises).\\
\texttt{employees.csv} & 600 rows & Skill, tenure, training and productivity; training is \emph{not} randomly assigned.\\\bottomrule
\end{tabularx}
\end{center}

\paragraph{The real datasets} All are free for teaching. Cite them in your work (full references in Appendix B).
\begin{center}\small
\begin{tabularx}{\linewidth}{@{}L{3.3cm}L{0.9cm}Y L{2.2cm}@{}}
\toprule
\textbf{Dataset} & \textbf{Week} & \textbf{What it is and why it is hard} & \textbf{Licence}\\\midrule
Online Retail II (UCI) & 1, 3, 4, 6 & Every transaction of a UK online gift wholesaler, Dec 2009--Dec 2011: 1\,067\,371 lines. Cancellations, write-offs, postage and fee codes mixed with sales, 243\,007 lines without a customer, a mixed-type invoice column, lumpy demand. & CC BY 4.0\\
Receipt phase of an environmental permit process (4TU) & 2 & Real event log from a Dutch municipality: 1\,434 cases, 8\,577 events, 116 variants, case deadlines. & 4TU terms of use\\
NAB machine temperature (Numenta) & 2 & Five-minute temperature of a large industrial machine with a planned shutdown and a catastrophic failure, labelled by experts. & Open (GitHub)\\
Hillstrom e-mail test & 7 & 64\,000 customers randomised to a men's e-mail, a women's e-mail or none; visits, conversions and spend over two weeks. & Free, public\\
Organ-donation panel & 7 & Quarterly donor-registration rates for 27 US states around California's 2011 active-choice rule. & In \texttt{causaldata}\\
Open-Meteo (optional) & 3 & Historical and forecast weather for any city, no API key. & Free, non-commercial\\\bottomrule
\end{tabularx}
\end{center}

\begin{insight}{Pre-course readiness check (self-study, about 3 hours)}
The course assumes you can already do the following without help. If you cannot, work through Appendix C before
Week 1.
\begin{enumerate}[nosep]
\item Load a CSV with pandas, parse dates, filter rows, group and aggregate, merge two tables, pivot.
\item Write a SQL query with \texttt{JOIN}, \texttt{GROUP BY}, \texttt{HAVING} and \texttt{ORDER BY}.
\item Fit and evaluate a scikit-learn regression model and explain overfitting.
\item Commit and push a notebook to GitHub.
\item State the four levels of analytics (descriptive, diagnostic, predictive, prescriptive) and give one
operations example of each. That is all the course will say about them.
\end{enumerate}
\end{insight}
